\chapter{The maximum principle}
\label{chap:maximum-principle}
\label{secmaxprin}
 Recall that the maximum principle for the heat equation says that if $h$ is a  solution to
the heat equation $$\frac{\partial h}{\partial t}=\Delta h$$ on a
compact manifold and if $h(x,0)\ge 0$ for all $x\in M$, then
$h(x,t)\ge 0$ for all $(x,t)$. In this chapter we discuss analogues
of this result for the scalar curvature, the Ricci curvature, and
the sectional curvature under Ricci flow. Of course, in all three
cases we are working with quasi-linear versions of the heat equation
so it is important to control the lower order (non-linear) terms and
in particular show that at zero curvature they have the appropriate
sign. Also, in the latter two cases we are working with tensors
rather than with scalars and hence we require a tensor version of
the maximum principle, which was established by Hamilton in
\cite{Hamilton4MPIC}.

As further applications of these results beyond just establishing
non-negativity, we indicate Hamilton's result that if the initial
conditions have positive Ricci curvature then the solution becomes
singular at finite time and as it does it becomes round (pinching to
round). We also give Hamilton's result showing that at points where
the scalar curvature is sufficiently large the curvature is pinched
toward positive. This result is crucial for understanding
singularity development. As a last application, we give Hamilton's
Harnack inequality for Ricci flows of non-negative curvature.


The maximum principle is used here in two different ways. The first
assumes non-negativity of something (e.g., a curvature) at time zero
and uses the maximum principle to establish non-negativity of this
quantity at all future times. The second assumes non-negativity of
something at all times and positivity at one point, and then uses
the maximum principle to establish positivity at all points and all
later times. In the latter application one compares the solution
with a solution to the linear heat equation where such a property is
known classically to hold.


\section{Maximum principle for scalar curvature}

Let us begin with the
easiest evolution equation, that for the scalar curvature, where the
argument uses only the (non-linear) version of the maximum
principle. This result is valid in all dimensions:


\begin{proposition}[Evolution bound for minimum scalar curvature]
\label{prop:scalar-curvature-min-evolution}
\uses{def:ricci-flow,def:ricci-curvature}\label{scalarevol}
Let $(M,g(t)),\ 0\le t<T$, be a Ricci flow with $M$ a compact
$n$-dimensional manifold. Denote by $R_\mathrm{min}(t)$ the minimum
value of the scalar curvature of $(M,g(t))$. Then:
\begin{itemize} \item $R_\mathrm{min}(t)$ is a non-decreasing function
of $t$.
\item If $R_\mathrm{min}(0)\ge 0$, then
$$R_\mathrm{min}(t)\ge R_\mathrm{min}(0)\left(\frac{1}{1-\frac{2t}{n}R_\mathrm{min}(0)}\right),$$ in particular,
$$T\le\frac{n}{2R_\mathrm{min}(0)}.$$
\item If $R_\mathrm{min}(0)<0$, then
$$R_\mathrm{min}(t)\ge
-\frac{n\bigl| R_\mathrm{min}(0)\bigr|}{2t\bigl|R_\mathrm{min}(0)\bigr|+n}.$$
\end{itemize}
\end{proposition}


\begin{proof}
\uses{lem:forward-difference-quotient,thm:curvature-evolution}
According to Equation~(\ref{Revol}), the evolution equation for $R$
is
$$\frac{\partial}{\partial t}R(x,t)=\Delta
R(x,t)+2|\Ric(x,t)|^2.$$ Since $M$ is compact, the function
$R_\mathrm{min}(t)$ is continuous but may not be $C^1$ at points where
the minimum of the scalar curvature is achieved at more than one
point.

The first thing to notice is the following:

\begin{claim}[Time derivative of R at its minimum]
\label{claim:scalar-min-derivative-bound}
\uses{def:ricci-curvature}\label{firstclaim}  If $R(x,t)=R_\mathrm{min}(t)$ then $(\partial R/\partial t)(x,t)\ge \frac{2}{n}R^2(x,t)$.
\end{claim}

\begin{proof}
This is immediate from  the evolution equation for $R$, the fact
that if $R(x,t)=R_\mathrm{min}(t)$, then $\Delta R(x,t)\ge 0$, and the
fact that $R$ is the trace of $\Ric$ which implies by the
Cauchy-Schwarz inequality that $|R|^2\le n|\Ric|^2$.
\end{proof}


Now it follows that:

\begin{claim}[Forward-difference bound for the minimum scalar curvature]
\label{claim:scalar-min-forward-difference}
\uses{claim:scalar-min-derivative-bound}\label{Rminprime}
$$\frac{d}{d t}(R_\mathrm{min}(t))\ge \frac{2}{n}R_\mathrm{min}^2(t),$$
where, at times $t$ where $R_\mathrm{min}(t)$ is not smooth, this
inequality is interpreted as an inequality for the forward
difference quotients.
\end{claim}


\begin{proof}
\uses{prop:forward-difference-maximum}
This is immediate from the first statement in
Proposition~\ref{fordiffmax}.
\end{proof}





If follows immediately from Claim~\ref{Rminprime} and
Lemma~\ref{fordiffquot} that $R_\mathrm{min}(t)$ is a non-decreasing
function of $t$. This establishes the first item and also the second
item in the case when $R_\mathrm{min}(0)=0$.

Suppose that $R_\mathrm{min}(0)\not= 0$.
 Consider the function
$$S(t)=\frac{-1}{R_\mathrm{min}(t)}-\frac{2t}{n}+\frac{1}{R_\mathrm{min}(0)}.$$
 Clearly, $S(0)=0$ and $S'(t)\ge 0$ (in the sense of forward difference quotients),
  so that by  Lemma~\ref{fordiffquot} we have $S(t)\ge 0$ for all $t$. This means that
\begin{equation}\label{Rminineq} \frac{1}{R_\mathrm{min}(t)}\le
\frac{1}{R_\mathrm{min}(0)}-\frac{2t}{n}\end{equation} provided that
$R_\mathrm{min}$ is not ever zero on the interval $[0,t]$. If $R_\mathrm{min}(0)>0$, then by the first item, $R_\mathrm{min}(t)>0$ for all $t$
for which the flow is defined, and the inequality in the second item
of the proposition is immediate from Equation~(\ref{Rminineq}). The
third inequality in the proposition also follows easily from
Equation~(\ref{Rminineq}) when $R_\mathrm{min}(t)<0$. But if $R_\mathrm{min}(t)\ge 0$, then the third item is obvious.
\end{proof}




\section{The maximum principle for tensors}

For the applications to the
Ricci curvature and the curvature tensor we need a version of the
maximum principle for tensors that is due to Hamilton; see
\cite{Hamilton4MPCO}.

Suppose that $V$ is a finite-dimensional real vector space and $Z\subset V$ is
a closed convex set. For each $z$ in the frontier of $ Z$ we define the {\em
tangent cone to} $Z$ at $z$, denoted $T_zZ$, to be the intersection of all
closed half-spaces $H$ of $V$ such that $z\in
\partial H$ and $Z\subset H$. For $z\in \mathrm{int}\, Z$ we define
$T_zZ=V$. Notice that $v\notin T_zZ$ if and only if there is a
affine linear function $\ell$ vanishing at $z$ non-positive on $Z$
and positive on $v$.

\begin{definition}[Vector field preserving a convex set]
\label{def:vector-field-preserves-convex-set}
Let $Z$ be a closed convex subset of a finite-dimensional real vector space
$V$. We say that a smooth vector field $\psi$ defined on an open neighborhood
$U$ of $Z$ in $V$ {\em preserves} $Z$ if for every $z\in Z$ we have $\psi(z)\in
T_zZ$.
\end{definition}

It is an easy exercise to show the following; see Lemma 4.1 on page
183 of \cite{Hamilton4MPCO}:

\begin{lemma}[Integral-curve criterion for preserving a convex set]
\label{lem:integral-curve-preserves-convex-set}
\uses{def:vector-field-preserves-convex-set}
Let $Z$ be a closed convex subset in a finite dimensional real
vector space $V$. Let $\psi$ be a smooth vector field defined on an
open neighborhood of $Z$ in $V$. Then $\psi$ preserves $Z$ if and
only if  every integral curve $\gamma\colon [0,a)\to V$ for $\psi$
with $\gamma(0)\in Z$ has $\gamma(t)\in Z$ for all $t\in [0,a)$.
Said more informally, $\psi$ preserves $Z$ if and only if every
integral curve for $\psi$ that starts in $Z$ remains in $Z$.
\end{lemma}




\subsection{The global
version}\label{4.3.3}

The maximum principle for tensors generalizes this to tensor flows
evolving by parabolic equations. First we introduce a generalization
of the notion of a vector field preserving a closed convex set to
the context of vector bundles.

\begin{definition}[Convex subbundle and fiberwise preserving vector fields]
\label{def:convex-subbundle-preserving-vector-field}
\uses{def:vector-field-preserves-convex-set}
 Let $\pi\colon {\mathcal V}\to M$ be a vector bundle and let
${\mathcal Z}\subset {\mathcal V}$ be  a closed subset. We say that
${\mathcal Z}$ is {\em convex} if  for every $x\in M$ the fiber
$Z_x$ of ${\mathcal Z}$ over $x$ is a convex subset of the vector
space fiber $V_x$ of ${\mathcal V}$ over $x$. Let $\psi$ be a
fiberwise vector field on an open neighborhood ${\mathcal U}$ of
${\mathcal Z}$ in ${\mathcal V}$. We say that $\psi$ {\em preserves}
${\mathcal Z}$ if for each $x\in M$ the restriction of $\psi$ to the
fiber $U_x$ of ${\mathcal U}$ over $x$ preserves $Z_x$.
\end{definition}

The following global version of the maximum principle for tensors is
Theorem 4.2 of  \cite{Hamilton4MPCO}.

\begin{theorem}[Maximum principle for tensors (global version)]
\label{thm:maximum-principle-tensors-global}
\uses{def:convex-subbundle-preserving-vector-field,thm:levi-civita-connection}\label{MP}{\bf (The maximum principle for tensors)}
Let $(M,g)$ be a compact Riemannian manifold. Let ${\mathcal V}\to M$ be a
tensor bundle and let ${\mathcal Z}\subset {\mathcal V}$ be a closed, convex
subset invariant under the parallel translation induced by the Levi-Civita
connection. Suppose that $\psi$ is a fiberwise vector field defined on an open
neighborhood of ${\mathcal Z}$ in ${\mathcal V}$ that preserves ${\mathcal Z}$.
Suppose that ${\mathcal T}(x,t),\ 0\le t\le T$, is a one-parameter family of
sections of ${\mathcal V}$ that evolves according to the parabolic equation
$$\frac{\partial {\mathcal T}}{\partial t}=\Delta {\mathcal T}
+\psi({\mathcal T}).$$ If ${\mathcal T}(x,0)$ is contained in
${\mathcal Z}$ for all $x\in M$, then ${\mathcal T}(x,t)$ is
contained in ${\mathcal Z}$ for all $x\in M$ and for all $0\le t\le
T$.
\end{theorem}


For a proof we refer the reader to  Theorem 4.3 and its proof (and
the related Theorem 4.2 and its proof) in  \cite{Hamilton4MPCO}.

There is a slight improvement of this result where the convex set ${\mathcal
Z}$ is allowed to vary with $t$. It is proved by the same argument; see Theorem
4.8 on page 101 of \cite{ChowKnopf}.


\begin{theorem}[Maximum principle for tensors with time-varying convex sets]
\label{thm:maximum-principle-tensors-varying-z}
\uses{def:convex-subbundle-preserving-vector-field,thm:levi-civita-connection}\label{varyingZ}
Let $(M,g)$ be a compact Riemannian manifold. Let ${\mathcal V}\to M$ be a
tensor bundle and let ${\mathcal Z}\subset {\mathcal V}\times [0,T]$ be a
closed subset with the property that for each $t\in [0,T]$ the time-slice
${\mathcal Z}(t)$ is a  convex subset of ${\mathcal V}\times\{t\}$ invariant
under the parallel translation induced by the Levi-Civita connection. Suppose
that $\psi$ is a fiberwise vector field defined on an open neighborhood of
${\mathcal Z}$ in ${\mathcal V}\times [0,T]$ that preserves the family
${\mathcal Z}(t)$ in the sense that any integral curve $\gamma(t),\ t_0\le t\le
t_1$, for $\psi$ with the property that $\gamma(t_0)\in {\mathcal Z}(t_0)$ has
$\gamma(t)\in {\mathcal Z}(t)$ for every $t\in [t_0,t_1]$. Suppose that
${\mathcal T}(x,t),\ 0\le t\le T$, is a one-parameter family of sections of
${\mathcal V}$ that evolves according to the parabolic equation
$$\frac{\partial {\mathcal T}}{\partial t}=\Delta {\mathcal T}+\psi({\mathcal T}).$$
If ${\mathcal T}(x,0)$ is contained in
${\mathcal Z}(0)$ for all $x\in M$, then ${\mathcal T}(x,t)$ is
contained in ${\mathcal Z}(t)$ for all $x\in M$ and for all $0\le
t\le T$.
\end{theorem}





\subsection{The local version}

Here is the local result. It is proved by the same argument as given
in the proof of Theorem 4.3 in \cite{Hamilton4MPCO}.

\begin{theorem}[Local maximum principle for tensors]
\label{thm:maximum-principle-tensors-local}
\uses{def:convex-subbundle-preserving-vector-field,thm:levi-civita-connection}\label{localMP}
Let $(M,g)$ be a Riemannian manifold. Let $\overline U\subset M$ be a compact,
smooth, connected, codimension-$0$ submanifold. Let ${\mathcal V}\to M$ be a
tensor bundle and let ${\mathcal Z}\subset {\mathcal V}$ be a closed, convex
subset. Suppose that $\psi$ is a fiberwise vector field defined on an open
neighborhood of ${\mathcal Z}$ in ${\mathcal V}$ preserving ${\mathcal Z}$.
Suppose that ${\mathcal Z}$ is invariant under the parallel translation induced
by the Levi-Civita connection. Suppose that ${\mathcal T}(x,t),\ 0\le t\le T$,
is a one-parameter family of sections of ${\mathcal V}$ that evolves according
to the parabolic equation
$$\frac{\partial {\mathcal T}}{\partial t}=\Delta {\mathcal T}
+\psi({\mathcal T}).$$ If ${\mathcal T}(x,0)$ is contained in
${\mathcal Z}$ for all $x\in \overline U$ and if ${\mathcal
T}(x,t)\in {\mathcal Z}$ for all $x\in \partial \overline U$ and all
$0\le t\le T$, then ${\mathcal T}(x,t)$ is contained in ${\mathcal
Z}$ for all $x\in \overline U$ and all $0\le t\le T$.
\end{theorem}



\section{Applications of the maximum principle}

Now let us give some applications of these results to Riemann and
Ricci curvature. In order to do this we first need to specialize the
above general maximum principles for tensors to the situation of the
curvature.

\subsection{Ricci flows with normalized initial conditions}

As we have already seen, the Ricci flow equation is invariant under multiplying
space and time by the same scale. This means that there can be no absolute
constants in the results about Ricci surgery. To break this gauge symmetry and
make the constants absolute we impose scale fixing (or rather scale bounding)
conditions on the initial metrics of the flows that we shall consider. The
following definition makes precise the exact conditions that we shall use.













\begin{definition}[Normalized initial conditions for Ricci flow]
\label{def:normalized-initial-conditions}
\uses{def:ricci-flow,def:riemann-curvature-tensor}\label{norminitcond}
We say that a that a Ricci flow $(M,g(t))$ has {\em normalized
initial conditions}
if $0$ is the initial time for the flow and if the compact
Riemannian manifold $(M^n,g(0))$ satisfies:
\begin{enumerate}
\item $|\Rm(x,0)|\le 1$ for all $x\in M$.
\item Let $\omega_n$ be the volume of the ball of radius $1$ in $n$ -dimensional Euclidean
space. Then $\Vol(B(x,0,r))\ge (\omega_n/2)r^n$ for any $p\in
M$ and any $r\le 1$.
\end{enumerate}
We also use the terminology $(M,g(0))$ {\em is normalized} to
indicate that it satisfies these two conditions.
\end{definition}

The evolution equation for the Riemann curvature and a standard maximum
principle argument show that if $(M,g(0))$ has an upper bound on the Riemann
curvature and a lower bound on the volume of balls of a fixed radius, then the
flow has Riemann curvature bounded above and volumes of balls bounded below on
a fixed time interval. Here is the result in the context of normalized initial
condition.



\begin{proposition}[Short-time curvature and volume bounds for normalized flows]
\label{prop:normalized-flow-short-time-bounds}
\uses{def:normalized-initial-conditions}\label{kappa0r0t0}
 There is $\kappa_0>0$  depending only on the dimension $n$
 such that the following holds.
Let $(M^n,g(t)),\ 0\le t\le T$, be a Ricci flow with bounded curvature, with
each $(M,g(t))$ being complete, and with normalized initial conditions. Then
$|\Rm(x,t)|\le 2$ for all $x\in M$ and all $t\in [0,\mathrm{min}(T,2^{-4})]$.
Furthermore, for any $t\in [0,\mathrm{min}(T,2^{-4})]$ and any $x\in M$ and any
$r\le 1$ we have $\Vol\,B(x,t,r)\ge \kappa_0r^n$.
\end{proposition}

\begin{proof}
\uses{thm:curvature-evolution}
The bound on the Riemann curvature follows directly from Lemma 6.1
on page 207 of \cite{ChowLuNi} and the definition of normalized
initial conditions. Once we know that the Riemann curvature is
bounded by $2$ on $[0,2^{-4}]$, there is an $0<r_0$ depending on $n$
such that for every $x\in M$ and every $r\le r_0$ we have
$B(x,0,r_0r)\subset B(x,t,r)\subset B(x,0,1)$. Also, from the bound
on the Riemann curvature and the evolution equation for volume given
in Equation~(\ref{Revol}), we see that there is $A<\infty$ such that
$\Vol_t\,(B(x,0,s))\ge A^{-1}\Vol_0\,(B(x,0,s))$. Putting
this together we see that
$$\Vol_t\,(B(x,t,r)\ge A^{-1}(\omega_n/2)r_0^nr^n.$$
This proves the result.
\end{proof}



\subsection{Extending flows}

There is one other consequence that will be important for us. For a
reference see \cite{ChowLuNi} Theorem 6.3 on page 208.


\begin{proposition}[Extension criterion for Ricci flows]
\label{prop:ricci-flow-extension-criterion}
\uses{def:ricci-flow}\label{flowextend}
Let $(M,g(t)),\ 0\le t<T<\infty$, be a Ricci flow with $M$ a compact manifold.
Then either the flow extends to an interval $[0,T')$ for some $T'>T$ or $|\Rm|$ is unbounded on $M\times[0,T)$.
\end{proposition}






\subsection{Non-negative curvature is preserved}

We need to consider the tensor versions of the maximum principle
when the tensor in question is the Riemann or Ricci curvature and
the evolution equation is that induced by the Ricci flow. This part
of the discussion is valid in dimension three only. We begin by
evaluating the expressions in Equation~(\ref{Mevol}) in the
$3$-dimensional case. Fix a symmetric bilinear form ${\mathcal S}$
on a $3$-dimensional real vector space $V$ with a positive definite
inner product. The inner product determines an identification of
$\wedge^2V$ with $V^*$. Hence, $\wedge^2{\mathcal S}^*$ is
identified with a symmetric automorphism of $V$,  denoted by
${\mathcal S}^\sharp$.

\begin{lemma}[Curvature-operator evolution ODE in dimension three]
\label{lem:curvature-operator-ode-dimension-three}
\uses{prop:curvature-operator-evolution,def:curvature-operator}\label{quad}
Let $(M,g)$ be a Riemannian $3$-manifold. Let ${\mathcal T}\in \Sym^2(\wedge^2T_x^*M)$ be the curvature operator written with
respect to the evolving frame as in Proposition~\ref{Mevol}. Then
the evolution equation given in Proposition~\ref{Mevol} is:
$$\frac{\partial {\mathcal T}}{\partial t}=\triangle {\mathcal T}+\psi({\mathcal T})$$
where
$$\psi({\mathcal T})={\mathcal T}^2+{\mathcal T}^\sharp.$$
In particular, in an orthonormal basis in which
$${\mathcal T}=\begin{pmatrix} \lambda & 0 & 0 \\ 0 & \mu & 0 \\ 0 &
0  & \nu\end{pmatrix}$$ with $\lambda\ge \mu\ge \nu$, the vector
field is given by
$$\psi({\mathcal T})={\mathcal T}^2+{\mathcal T}^\sharp=\begin{pmatrix}
\lambda^2+\mu\nu & 0 & 0 \\ 0 & \mu^2+\lambda\nu & 0 \\ 0 & 0 &
\nu^2+\lambda\mu\end{pmatrix}.$$
\end{lemma}

\begin{corollary}[Non-negative curvature operator is preserved]
\label{cor:nonnegative-curvature-operator-preserved}
\uses{def:curvature-operator}\label{rm>0}
Let $(M,g(t)),\ 0\le t\le T$, be a Ricci flow with $M$ a compact,
connected $3$-manifold. Suppose that  $\Rm(x,0)\ge 0$ for all
$x\in M$. Then $\Rm(x,t)\ge 0$ for all $x\in M$ and all $t\in
[0,T]$.
\end{corollary}

\begin{proof}
\uses{thm:maximum-principle-tensors-global,lem:curvature-operator-ode-dimension-three}
Let $\nu_x\colon \Sym^2(\wedge^2T^*_xM)\to \Ar$ associate to
each endomorphism its smallest eigenvalue. Then $\nu_x({\mathcal
T})$ is the minimum over all lines in $\wedge^2T_xM$ of the trace of
the restriction of ${\mathcal T}$ to that line. As a minimum of
linear functions, $\nu_x$ is a convex function. In particular,
$Z_x=\nu_x^{-1}([0,\infty))$ is a convex subset. We let ${\mathcal
Z}$ be the union over all $x$ of $Z_x$. Clearly, ${\mathcal Z}$ is a
closed convex subset of the tensor bundle. Since parallel
translation is orthogonal, ${\mathcal Z}$ is invariant under
parallel translation. The expressions in Lemma~\ref{quad} show that
if ${\mathcal T}$ is an endomorphism of $\wedge^2T^*_xM$ with
$\nu({\mathcal T})\ge 0$, then the symmetric matrix $\psi({\mathcal
T})$ is non-negative. This implies that $\nu_x$ is non-decreasing in
the direction $\psi({\mathcal T})$ at the point ${\mathcal T}$. That
is to say, for each $x\in M$, the vector field $\psi({\mathcal T})$
preserves the set $\{\nu_x^{-1}([c,\infty))\}$ for any $c\ge 0$. The
hypothesis that $\Rm(x,0)\ge 0$ means that $\Rm(x,0)\in
{\mathcal Z}$ for all $x\in M$. Applying Theorem~\ref{MP} proves the
result.
\end{proof}


\begin{corollary}[Non-negative Ricci curvature is preserved]
\label{cor:nonnegative-ricci-preserved}
\uses{def:ricci-curvature}\label{Ric>0}
Suppose that $(M,g(t)),\ 0\le t\le T$, is a Ricci flow with $M$ a
compact, connected $3$-manifold with $\Ric(x,0)\ge 0$ for all
$x\in M$. Then $\Ric(x,t)\ge 0$ for all $t>0$.
\end{corollary}


\begin{proof}
\uses{thm:maximum-principle-tensors-global,lem:curvature-operator-ode-dimension-three}
The statement that $\Ric(x,t)\ge 0$ is equivalent to the statement that
for every two-plane in $\wedge^2 T_xM$ the trace of the Riemann curvature
operator on this plane is $\ge 0$. For ${\mathcal T}\in \Sym^2(\wedge^2T_x^*M)$, we define $s({\mathcal T})$ as the minimum over all
two-planes $P$ in $\wedge^2TM$ of the trace of ${\mathcal T}$ on $P$. The
restriction $s_x$ of $s$ to the fiber over $x$ is the minimum of a collection
of linear functions and hence is convex. Thus, the subset ${\mathcal
S}=s^{-1}([0,\infty))$ is convex. Clearly, $s$ is preserved by orthogonal
isomorphisms, so ${\mathcal S}$ is invariant under parallel translation. Let
$\lambda\ge\mu\ge\nu$ be the eigenvalues of ${\mathcal T}$. According to
Lemma~\ref{quad} the derivative of $s_x$ at ${\mathcal T}$ in the
$\psi({\mathcal T})$-direction is
$(\mu^2+\lambda\nu)+(\nu^2+\lambda\mu)=(\mu^2+\nu^2)+\lambda(\mu+\nu)$. The
condition that $s({\mathcal T})\ge 0$, is the condition that $\nu+\mu\ge 0$,
and hence $\mu\ge 0$, implying that $\lambda\ge 0$. Thus, if $s({\mathcal
T})\ge 0$, it is also the case that the derivatifve of $s_x$ in the
$\psi({\mathcal T})$-direction is non-negative. This implies that $\psi$
preserves ${\mathcal S}$. Applying Theorem~\ref{MP} gives the result.
\end{proof}


\section{The strong maximum principle for curvature}

First let us state the strong maximum principle for the heat
equation.% see \cite{??}.

\begin{theorem}[Strong maximum principle for the heat equation]
\label{thm:strong-maximum-principle-heat}\label{SMPheat}
Let $\overline U$ be a compact, connected manifold, possibly with boundary. Let
$h(x,t),\ 0\le t\le T$, be a solution to the heat equation
$$\frac{\partial h(x,t)}{\partial t}=\Delta h(x,t).$$
 Suppose that  $h$ has Dirichlet boundary conditions in the sense
that $h(x,t)=0$ for all $(x,t)\in\partial \overline U\times [0,T]$.
If $h(x,0)\ge 0$ for all $x\in \overline U$, then $h(x,t)\ge 0$ for
all $(x,t)\in \overline U\times [0,T]$. If, in addition, there is
$y\in \overline U$ with $h(y,0)>0$, then $h(x,t)>0$ for all
$(x,t)\in \mathrm{int}(\overline U)\times(0,T]$.
\end{theorem}

We shall use this strong maximum principle to establish an analogous
result for the curvature tensors. The hypotheses are in some ways
more restrictive -- they are set up to apply to the Riemann and
Ricci curvature.

\begin{proposition}[Strong maximum principle for curvature tensors]
\label{prop:strong-maximum-principle-curvature}
\uses{def:convex-subbundle-preserving-vector-field}\label{SMPcurv}
Let $(M,g)$ be a Riemannian manifold and let ${\mathcal V}$ be a
tensor bundle. Suppose that $\overline U$ is a compact, connected,
smooth codimension-$0$ submanifold of $M$. Consider a one-parameter
family of sections ${\mathcal T}(x,t),\ 0\le t\le T$, of ${\mathcal
V}$. Suppose that ${\mathcal T}$ evolves according to the equation
$$\frac{\partial {\mathcal T}}{\partial t}=\Delta {\mathcal T}+\psi({\mathcal T})$$
for some smooth, fiberwise vector field $\psi({\mathcal T})$
defined on ${\mathcal V}$. Suppose that $s\colon {\mathcal V}\to
\Ar$ is a function satisfying the following properties:
\begin{enumerate}
\item For each $x\in M$ the restriction $s_x$ to the fiber $V_x$ of
${\mathcal V}$ over $x$ is a  convex function.
\item For any $A$ satisfying $s_x(A)\ge 0$ the vector $\psi(A)$ is contained
in the tangent cone of the convex set $\{y|s_x(y)\ge s_x(A)\}$ at the
point $A$.
\item $s$ is invariant under parallel translation.
\end{enumerate}
Suppose that $s({\mathcal T}(x,0))\ge 0$ for all $x\in \overline U$
and that $s({\mathcal T}(x,t))\ge 0$ for all $x\in \partial\overline
U$ and all $t\in [0,T]$. Suppose also that there is $x_0\in \mathrm{int}(\overline U)$ with $s({\mathcal T}(x_0,0))>0$. Then
$s({\mathcal T}(x,t))> 0$ for all $(x,t)\in \mathrm{int}(\overline
U)\times(0,T]$.
\end{proposition}


\begin{proof}
\uses{thm:strong-maximum-principle-heat,thm:maximum-principle-tensors-local}
Let $h\colon \overline U\times\{0\}\to \Ar$ be a smooth function
with $h(x,0)=0$ for all $x\in \partial \overline U$ and with
$s({\mathcal T}(x,0))\ge h(x,0)\ge 0$ for all $x\in \overline U$. We
choose $h$ so that $h(x_0,0)>0$. Let $h(x,t),\ 0\le t<\infty$,  be
the solution to the heat equation on $\overline U$
$$\frac{\partial h}{\partial t}=\Delta h$$
with Dirichlet boundary conditions $h(x,t)=0$ for all $x\in \partial
\overline U$ and all $t\ge 0$ and with the given initial conditions.


Consider the tensor bundle ${\mathcal V}\oplus \Ar$ over $M$. We
define
$$Z_x=\left\{({\mathcal T},h)\in V_x\oplus
\Ar\bigl|\bigr. s_x({\mathcal T})\ge h\ge 0\right\}.$$ The union
over all $x\in M$ of the $Z_x$ defines a closed convex subset
${\mathcal Z}\subset {\mathcal V}\oplus \Ar$ which is invariant
under parallel translation since $s$ is. We consider the family of
sections $({\mathcal T}(x,t),h(x,t)),\ 0\le t\le T$, of ${\mathcal
V}\oplus \Ar$. These evolve by
$$\frac{d\left({\mathcal T}(x,t),h(x,t)\right)}{dt}=\left(\Delta {\mathcal T}(x,t),\Delta
h(x,t)\right)+\widetilde \psi\left({\mathcal T}(x,t),h(x,t)\right)$$
where $\widetilde \psi({\mathcal T},h)=\left(\psi({\mathcal
T}),0\right)$. Clearly, by our hypotheses, the vector field
$\widetilde \psi$ preserves the convex set ${\mathcal Z}$. Applying
the local version of the maximum principle (Theorem~\ref{localMP}),
we conclude that ${\mathcal T}(x,t)\ge h(x,t)$ for all $(x,t)\in
\overline U\times [0,T]$.

The result then follows immediately from Theorem~\ref{SMPheat}.
\end{proof}




\subsection{Applications of the strong maximum principle}

We have the following applications of the strong maximum principle.




\begin{theorem}[Vanishing scalar curvature forces a flat metric]
\label{thm:vanishing-scalar-curvature-forces-flat}\label{firststrongmax}
Let $(U,g(t)),\ 0\le t\le T$, be a $3$-dimensional Ricci flow with
non-negative sectional curvature with $U$ connected but not
necessarily complete and with $T>0$. If $R(p,T)=0$ for some $p\in
U$, then $(U,g(t))$ is flat for every $t\in [0,T]$.
\end{theorem}

\begin{proof}
\uses{thm:curvature-evolution}
We suppose that there is $p\in U$ with $R(p,T)=0$. Since all the
metrics in the flow are of non-negative sectional curvature, if the
flow does not consist entirely of flat manifolds then there is
$(q,t)\in U\times [0,T]$ with $R(q,t)>0$. Clearly, by continuity, we
can assume $t<T$. By restricting to the time interval $[t,T]$ and
shifting by $-t$ we can arrange that $t=0$. Let $V$ be a compact,
connected smooth submanifold with boundary whose interior contains
$q$ and $p$. Let $h(y,0)$ be a smooth non-negative function with
support in $V$, positive at $q$, such that $R(y,0)\ge h(y,0)$ for
all $y\in V$. Let $h(y,t)$ be the solution to the heat equation on
$V\times [0,T]$ that vanishes on $\partial V$. Of course, $h(y,T)>0$
for all $y\in \mathrm{int}(V)$. Also, from Equation~(\ref{Revol}) we
have
$$\frac{\partial}{\partial t}(R-h)=\triangle (R-h)+2|\Ric|^2,$$
so that $(R-h)(y,0)\ge 0$ on $(V\times \{0\})\cup (\partial V\times [0,T])$. It
follows from the  maximum principle that $(R-h)\ge 0$ on all of $V\times
[0,T]$. In particular, $R(p,T)\ge h(p,T)>0$. This is a contradiction,
establishing the theorem.
\end{proof}

\begin{corollary}[Local product splitting at a Ricci null eigenvalue]
\label{cor:local-product-splitting-ricci-null-eigenvalue}
\uses{def:curvature-operator}\label{localprod}
Fix $T>0$. Suppose that $(U,g(t)),\ 0\le t\le T$, is a  Ricci flow
such that for each $t$, the Riemannian manifold  $(U,g(t))$ is a
(not necessarily complete) connected, $3$-manifold of non-negative
sectional curvature. Suppose that $(U,g(0))$ is not flat and that
for some $p\in M$ the Ricci curvature at $(p,T)$ has a zero
eigenvalue. Then for each $t\in (0,T]$ the Riemannian manifold
$(U,g(t))$ splits locally as a product of a surface of positive
curvature and a line, and under this local splitting the flow is
locally the product of a Ricci flow on the surface and the trivial
flow on the line.
\end{corollary}

\begin{proof}
\uses{cor:nonnegative-ricci-preserved,thm:curvature-evolution,prop:strong-maximum-principle-curvature,thm:vanishing-scalar-curvature-forces-flat}
First notice that it follows from Theorem~\ref{firststrongmax} that
because $(U,g(0))$ is not flat, we have $R(y,t)>0$ for every
$(y,t)\in U\times (0,T]$.

 We consider the function $s$ on
$\Sym^2(\wedge^2T^*_yU)$ that associates to each endomorphism
the sum of the smallest $2$ eigenvalues. Then $s_y$ is the minimum
of the traces on $2$-dimensional subsets  in $\wedge^2T_yU$. Thus,
$s$ is a convex function, and the subset ${\mathcal
S}=s^{-1}([0,\infty))$ is a convex subset. Clearly, this subset is
invariant under parallel translation. By the computations in the
proof of Corollary~\ref{Ric>0} it is invariant under the vector
field $\psi({\mathcal T})$. The hypothesis of the corollary tells us
that $s(p,T)=0$. Suppose that $s(q,t)>0$ for some $(q,t)\in U\times
[0,T]$. Of course, by continuity we can take $t< T$. Shift the time
parameter so that $t=0$, and fix a compact connected,
codimension-$0$ submanifold $V$ containing $p,q$ in its interior.
Then by Theorem~\ref{SMPcurv} $s(y,T)>0$ for all $y\in \mathrm{int}(V)$
and in particular $s(p,T)>0$. This is a contradiction, and we
conclude that $s(q,t)=0$ for all $(q,t)\in U\times [0,T]$.

Since we have already established that each $R(y,t)>0$ for all $(y,t)\in
U\times (0,T]$, so that $\Rm(y,t)$ is not identically zero, this means
that for all $y\in U$ and all $t\in (0,T]$ that the null space of the  operator
$\Rm(y,t)$ is a $2$-dimensional subspace of $\wedge^2T_yU$. This
$2$-dimensional subspace is dual to a line in $T_xM$. Thus, we have a
one-dimensional distribution (a line bundle in the tangent bundle) ${\mathcal
D}$ in $U\times (0,T]$ with the property that the sectional curvature $\Rm(y,t)$ vanishes on any $2$-plane containing the line ${\mathcal D}(y,t)$.
The fact that the sectional curvature of $g(t)$ vanishes on all two-planes in
$T_yM$ containing ${\mathcal D}(y,t)$ means that its eigenvalues are
$\{\lambda,0,0\}$ where $\lambda>0$ is the sectional curvature of the
$g(t)$-orthogonal $2$-plane to ${\mathcal D}(y,t)$. Hence ${\mathcal
R}(V(y,t),\cdot,\cdot,\cdot)=0$.

Locally in space and time, there is a unique (up to sign) vector field $V(y,t)$
that generates ${\mathcal D}$ and satisfies $|V(y,t)|^2_{g(t)}=1$. We wish to
show that this local vector field is invariant under parallel translation and
time translation; cf. Lemma 8.2 in \cite{Hamilton4MPCO}. Fix a point $x\in M$,
a direction $X$ at $x$, and a time $t$. Let $\tilde V(y,t)$ be a parallel
extension of $V(x,t)$ along a curve $C$ passing through $x$ in the
$X$-direction, and let $\tilde W(y,t)$ be an arbitrary parallel vector field
along $C$. Since the sectional curvature is non-negative, we have ${\mathcal
R}(\tilde V,\tilde W,\tilde V,\tilde W)(y)\ge 0$ for all $y\in C$; furthermore,
this expression vanishes at $x$. Hence, its first variation vanishes at $x$.
That is to say
$$\nabla \left({\mathcal R}(\tilde V,\tilde W,\tilde V,\tilde W)\right)(x,t)=(\nabla{\mathcal R})(\tilde
V,\tilde W,\tilde V,\tilde W)$$ vanishes at $(x,t)$. Since this is true for all
$\tilde W$, it follows that the null space of the quadratic form
$\nabla{\mathcal R}(x,t)$ contains the null space of ${\mathcal R}(x,t)$, and
thus $$(\nabla {\mathcal R})(V(x,t),\cdot,\cdot,\cdot)=0.$$
 Now let us consider three parallel
vector fields $\tilde W_1,\tilde W_2,$ and $\tilde W_3$ along $C$.
We compute $0=\nabla_X\left({\mathcal R}( V(y,t),\tilde
W_1(y,t),\tilde W_2(y,t),\tilde W_3(y,t))\right)$. (Notice that
while the $\tilde W_i$ are parallel along $C$, $V(y,t)$ is defined
to be the vector field spanning ${\mathcal D}(y,t)$ rather than a
parallel extension of $V(x,t)$.) Given the above result we find that
$$0=2{\mathcal R}(\nabla_X V(x,t),\tilde W_1(x,t),\tilde W_2(x,t),
\tilde W_3(x,t)).$$ Since this is true for all triples of vector
fields $\tilde W_i(x,t)$, it follows that $\nabla_XV(x,t)$ is a real
multiple of $V(x,t)$. But since $|V(y,t)|^2_{g(t)}=1$, we see that
$\nabla_XV(x,t)$ is orthogonal to $V(x,t)$. We conclude that
$\nabla_XV(x,t)=0$. Since $x$ and $X$ are general, this shows that
the local vector field $V(x,t)$ is invariant under the parallel
translation associated to the metric $g(t)$.

It follows that locally $(M,g(t))$ is a Riemannian product of a surface of
positive curvature with a line. Under this product decomposition, the curvature
is the pullback of the curvature of the surface. Hence, by
Equation~(\ref{Rmevol}), under Ricci flow on the $3$-manifold, the time
derivative of the curvature at time $t$ also decomposes as the pullback of the
time derivative of the curvature of the surface under Ricci flow on the
surface. In particular, $(\partial {\mathcal R}/\partial
t)(V,\cdot,\cdot,\cdot)=0$. It now follows easily that $\partial
V(x,t)/\partial t=0$.

This completes the proof that the unit vector field in the direction ${\mathcal
D}(x,t)$ is invariant under parallel translation and under time translation.
Thus, there is a local Riemannian splitting of the $3$-manifold into a surface
and a line, and this splitting is invariant under the Ricci flow. This
completes the proof of the corollary.
\end{proof}

In the complete case, this local product decomposition globalizes in
some cover; see Lemma 9.1 in \cite{Hamilton4MPCO}.


\begin{corollary}[Global product splitting on a cover]
\label{cor:global-product-splitting-cover}
\uses{cor:local-product-splitting-ricci-null-eigenvalue}\label{nullspace}
Suppose that $(M,g(t)),\ 0\le t\le T$, is a Ricci flow of complete, connected
Riemannian $3$-manifolds with $\Rm(x,t)\ge 0$ for all $(x,t)$ and with
$T>0$. Suppose that $(M,g(0))$ is not flat and that for some $x\in M$  the
endomorphism $\Rm(x,T)$ has a zero eigenvalue. Then $M$ has a cover
$\widetilde M$ such that, denoting the induced family of metrics on this cover
by $\tilde g(t)$, we have that $(\widetilde M,\tilde g(t))$ splits as a product
$$(N,h(t))\times (\Ar,ds^2)$$
where $(N,h(t))$ is a surface of positive curvature for all $0<t\le
T$. The Ricci flow is a product of the Ricci flow $(N,h(t)),\ 0\le
t\le T$, with the trivial flow on $\Ar$.
\end{corollary}


\begin{remark}[Classification of covers for the product splitting]
\label{rem:product-splitting-cover-classification}
\uses{cor:global-product-splitting-cover}
Notice that there are only four possibilities for the cover required
by the corollary. It can be trivial, or a normal $\Zee$-cover or it
can be a two-sheeted cover or a normal infinite dihedral group
cover. In the first two cases, there is a unit vector field on $M$
parallel under $g(t)$ for all $t$ spanning the null direction of
$\Ric$. In the last two cases, there is no such vector field,
only a non-orientable line field.
\end{remark}


 Let $(N,g)$ be a Riemannian manifold. Recall from Definition~\ref{conedefn}
  that the open cone
on $(N,g)$ is the space $N\times (0,\infty)$ with the Riemannian
metric $\tilde g(x,s)=s^2g(x)+ds^2$. An extremely important result
for us is that open pieces in non-flat cones cannot arise as the
result of Ricci flow with non-negative curvature.



\begin{proposition}[Non-flat cones do not arise from non-negatively curved Ricci flow]
\label{prop:no-nonflat-cones-from-ricci-flow}
\uses{def:cone}\label{nocones}
Suppose that $(U,g(t)),\ 0\le t\le T$, is a $3$-dimensional Ricci
flow with non-negative sectional curvature, with $U$ being connected
but not necessarily complete and $T>0$. Suppose that $(U,g(T))$ is
isometric to a non-empty open subset of a cone over a Riemannian
manifold. Then $(U,g(t))$ is flat for every $t\in [0,T]$.
\end{proposition}


\begin{proof}
\uses{prop:cone-curvature,thm:vanishing-scalar-curvature-forces-flat,cor:local-product-splitting-ricci-null-eigenvalue}
If $(U,g(T))$ is flat, then by Theorem~\ref{firststrongmax}
 for every $t\in [0,T]$ the
Riemannian manifold $(U,g(t))$ is flat.

We must rule out the possibility that $(U,g(T))$ is  non-flat. Suppose that
$(U,g(T))$ is an open subset in a non-flat cone. According to
Proposition~\ref{conecurv},  for each $x\in U$ the Riemann curvature tensor of
$(U,g(T))$ at $x$ has a $2$-dimensional null space in $\wedge^2T_xU$. Since we
are assuming that $(U,g(T))$ is not flat, the third eigenvalue of the Riemann
curvature tensor is not identically zero. Restricting to a smaller open subset
if necessary, we can assume that the third eigenvalue is never zero. By the
computations in Proposition~\ref{conecurv} the non-zero eigenvalue is not
constant, and in fact it scales by $s^{-2}$ in the terminology of that
proposition, as we move along the cone lines. Of course, the $2$-dimensional
null-space for the Riemann curvature tensor at each point is equivalent to a
line field in the tangent bundle of the manifold. Clearly, that line field is
the line field along the cone lines. Corollary~\ref{localprod} says that since
the Riemann curvature of $(U,g(T))$ has a $2$-dimensional null-space in
$\wedge^2T_xU$ at every point $x\in U$, the Riemannian manifold $(U,g(T))$
locally splits as a Riemannian product of a line with a surface of positive
curvature, and the $2$-dimensional null-space for the Riemannian curvature
tensor is equivalent to the line field in the direction of the second factor.
Along these lines the non-zero eigenvalue of the curvature is constant. This is
a contradiction and establishes the result.
\end{proof}

Lastly, we have Hamilton's result (Theorem 15.1 in
\cite{Hamilton3MPRC}) that compact $3$-manifolds of non-negative
Ricci curvature become round under Ricci flow:

\begin{theorem}[Compact non-negative Ricci flows become round]
\label{thm:nonnegative-ricci-flow-becomes-round}
\uses{cor:nonnegative-ricci-preserved}\label{flowtoround}
Suppose that $(M,g(t)),\ 0\le t<T$, is a Ricci flow with $M$ being a
compact $3$-dimensional manifold. If $\Ric(x,0)\ge 0$ for all
$x\in M$, then either $\Ric(x,t)>0$ for all $(x,t)\in M\times
(0,T)$ or $\Ric(x,t)=0$ for all $(x,t)\in M\times[0,T)$.
Suppose that $\Ric(x,t)>0$ for some $(x,t)$ and that the flow
is maximal in the sense that there is no $T'>T$ and an extension of
the given flow to a  flow defined on the time interval $[0,T')$. For
each $(x,t)$, let $\lambda(x,t)$, resp. $\nu(x,t)$, denote the
largest, resp. smallest, eigenvalue of $\Rm(x,t)$ on
$\wedge^2T_xM$. Then as $t$ tends to $T$ the Riemannian manifolds
$(M,g(t))$ are becoming round in the sense that
$$\mathrm{lim}_{t\rightarrow T}\frac{\mathrm{max}_{x\in M}\lambda(x,t)}{\mathrm{min}_{x\in M}\nu(x,t)}=1.$$ Furthermore, for any $x\in M$ the
largest eigenvalue $\lambda(x,t)$ tends to $\infty$ as $t$ tends to
$T$, and rescaling $(M,g(t))$ by $\lambda(x,t)$ produces a family of
Riemannian manifolds converging smoothly as $t$ goes to $T$ to a
compact round manifold. In particular, the underlying smooth
manifold supports a Riemannian metric of constant positive curvature
so that the manifold is diffeomorphic to a $3$-dimensional spherical
space-form.
\end{theorem}

Hamilton's proof in \cite{Hamilton3MPRC} uses the maximum principle
and Shi's derivative estimates.






\subsection{Solitons of positive curvature}

One nice application of this pinching result is the following
theorem.

\begin{theorem}[Positive Ricci solitons are round spherical space forms]
\label{thm:positive-ricci-soliton-round}\label{sphsf}
Let $(M,g)$ be a compact $3$-dimensional soliton of positive Ricci
curvature. Then $(M,g)$ is round. In particular, $(M,g)$ is the
quotient of $S^3$ with a round metric by a finite subgroup of $O(4)$
acting freely; that is to say, $M$ is a $3$-dimensional spherical
space-form.
\end{theorem}

\begin{proof}
\uses{thm:nonnegative-ricci-flow-becomes-round}
Let $(M,g(t)),\ 0\le t<T$, be the maximal Ricci flow with initial
manifold $(M,g)$. Since $\Ric(x,0)>0$ for all $x\in M$, it
follows from Theorem~\ref{flowtoround} that $T<\infty$ and that as
$t$ tends to $T$ the metrics $g(t)$ converge smoothly to a round
metric. Since all the manifolds $(M,g(t))$ are isometric up to
diffeomorphism and a constant conformal factor, this implies that
all the $g(t)$ are of constant positive curvature.

The last statement is a standard consequence of the fact that the
manifold has constant positive curvature.
\end{proof}


\begin{remark}[Extending roundness to shrinking solitons]
\label{rem:roundness-shrinking-soliton-preview}
\uses{thm:positive-ricci-soliton-round,def:gradient-shrinking-soliton}
After we give a stronger pinching result in the next section, we
shall improve this result, replacing the positive Ricci curvature
assumption by the {\em a priori} weaker assumption that the soliton
is a shrinking soliton.
\end{remark}




\section{Pinching toward positive
curvature}

As the last application of the maximum principle for tensors we give
a theorem due to R. Hamilton (Theorem 4.1 in \cite{HamiltonNSRF3M})
and T. Ivey \cite{Ivey} which shows that, in dimension three, as the
scalar curvature gets large, the sectional curvatures pinch toward
the positive. Of course, if the sectional curvatures are
non-negative, then the results in the previous section apply. Here,
we are considering the case when the sectional curvature is not
everywhere positive. The pinching result says roughly the following:
At points where the Riemann curvature tensor has a negative
eigenvalue, the smallest (thus negative) eigenvalue of the Riemann
curvature tensor divided by the largest eigenvalue limits to zero as
the scalar curvature grows. This result is central in the analysis
of singularity development in finite time for a $3$-dimensional
Ricci flow.

\begin{theorem}[Pinching toward positive curvature]
\label{thm:pinching-toward-positive-curvature}
\uses{def:ricci-flow,def:curvature-operator}\label{pinch}{\bf (Pinching toward positive curvature)}
Let $(M,g(t)),\ 0\le t<T$, be a Ricci flow with $M$ a compact
$3$-manifold. Assume at for every $x\in M$, the eigenvalues,
$\lambda(x,0)\geq\mu(x,0)\geq\nu(x,0)$, of $\Rm(x,t)$ are all
at least $-1$. Set $X(x,t)=\mathrm{max}(-\nu(x,t),0)$. Then we have:
\begin{enumerate}
\item $R(x,t)\ge \frac{-6}{4t+1}$, and
\item for all $(x,t)$ for which $0<X(x,t)$
$$R(x,t)\geq 2X(x,t)\left(\mathrm{log}X(x,t)+\mathrm{log}(1+t)-3\right).$$
\end{enumerate}
\end{theorem}

For any fixed $t$, the limit as $X$ goes to $0$ from above of
$X(\mathrm{log}(X)+\mathrm{log}(1+t)-3)$ is zero, so that it is natural to
interpret this expression to be zero when $X=0$. Of course, when
$X(x,t)=0$ all the eigenvalues of $\Rm(x,t)$ are non-negative
so that $R(x,t)\ge 0$ as well. Thus, with this interpretation of the
expression in Part 2 of the theorem, it remains valid even when
$X(x,t)=0$.

\begin{remark}[Consequence: negative eigenvalues are small relative to scalar curvature]
\label{rem:pinching-negative-eigenvalues-small}
\uses{thm:pinching-toward-positive-curvature}
This theorem tells us, among other things, that as the scalar
curvature goes to infinity  then  absolute values of all the
negative eigenvalues (if any) of $\Rm$ are  arbitrarily small
with respect to the scalar curvature.
\end{remark}


We proof we give below follows Hamilton's original proof in
\cite{HamiltonNSRF3M} very closely.

\begin{proof}
\uses{prop:scalar-curvature-min-evolution,thm:maximum-principle-tensors-varying-z}
First note that by Proposition~\ref{scalarevol}, if $R_\mathrm{min}(0)\ge 0$, then the same is true for $R_\mathrm{min}(t)$ for every
$t>0$ and thus the first inequality stated in the theorem is clearly
true. If $R_\mathrm{min}(0)<0$, the first inequality stated in the
theorem follows easily from the last inequality in
Proposition~\ref{scalarevol}.

We turn now to the second inequality in the statement of the
theorem. Consider the tensor bundle ${\mathcal V}=\Sym^2(\wedge^2T^*M)$. Then the curvature operator written in the
evolving frame, ${\mathcal T}(x,t)$, is a one-parameter family of
smooth sections of this bundle, evolving by
$$\frac{\partial{\mathcal T}}{\partial t}=
\Delta{\mathcal T}+\psi({\mathcal T}).$$ We consider two subsets of ${\mathcal
V}$. There are two solutions to $x(\mathrm{log}(x)+(\mathrm{log}(1+t)-3)=-3/(1+t)$.
One is $x=1/(1+t)$; let $\xi(t)>1/(1+t)$ be the other. We set $S({\mathcal
T})=\tr({\mathcal T})$, so that $R=2S$, and we set $X({\mathcal T})=\mathrm{max}(-\nu({\mathcal T}),0)$. Define
\begin{eqnarray*}
 {\mathcal Z}_1(t) & = & \{{\mathcal T}\in {\mathcal V}\bigl|\bigr.
 S({\mathcal T})\geq-\frac{3}{(1+t)}\} \\
{\mathcal Z}_2(t) & = & \{ {\mathcal T}\in {\mathcal
V}\bigl|\bigr.S({\mathcal T})\geq f_t(X({\mathcal T})),\ \ \
\text{if }\ \ X( {\mathcal T})\geq\xi(t)\},
\end{eqnarray*}
where $f_t(x)=x(\mathrm{log}x+\mathrm{log}(1+t)-3)$. Then we define
$${\mathcal Z}(t)={\mathcal Z}_1(t)\cap {\mathcal Z}_2(t).$$




\begin{claim}[Convexity of the pinching region]
\label{claim:pinching-region-convex}
For each $x\in M$ and each $t\ge 0$, the fiber $Z(x,t)$ of
${\mathcal Z}(t)$ over $x$  is a convex subset of $\Sym^2(\wedge^2T^*M)$.
\end{claim}


\begin{proof}
First consider the function $f_t(x)=x(\mathrm{log}(x)+\mathrm{log}(1+t)-3)$ on the interval $[\xi(t),\infty)$. Direct computation
shows that $f'(x)>0$ and $f''(x)>0$ on this interval. Hence, for
every $t\ge 0$ the region ${\mathcal C}(t)$ in the $S$-$X$ plane
 defined by $S\ge -3/(1+t)$
and $S\ge f_t(X)$ when $X\ge \xi(t)$ is convex and has the property
that if $(S,X)\in {\mathcal C}(t)$ then so is $(S,X')$ for all
$X'\le X$.. By definition an element
${\mathcal T}\in {\mathcal V}$ is contained in ${\mathcal Z}(t)$ if
and only if $(S({\mathcal T}),X({\mathcal T})\in {\mathcal C}(t)$.
Now fix $t\ge 0$ and suppose that ${\mathcal T}_1$ and ${\mathcal
T}_2$ are elements of $\Sym^2(\wedge^2T^*M_x)$ such that
setting $S_i=\tr({\mathcal T}_i)$ and $X_i=X({\mathcal T}_i)$
we have $(S_i,X_i)\in {\mathcal C}(t)$ for $i=1,2$. Then we consider
${\mathcal T}=s{\mathcal T}_1+(1-s){\mathcal T}_2$ for some $s\in
[0,1]$. Let $S=\tr({\mathcal T})$ and $X=X({\mathcal T})$.
Since ${\mathcal C}(t)$ is convex, we know that
$(sS_1+(1-s)S_2,sX_1+(1-s)X_2)\in{\mathcal C}(t)$, so that
${\mathcal T}\in {\mathcal Z}(t)$.
 Clearly,
$S=sS_1+(1-s)S_2$, so that we conclude that $(S,(sX_1+(1-s)X_2))\in{\mathcal
C}(t)$.  But since $\nu$ is a convex function, $X$ is a concave function, i.e.,
$X\le sX_1+(1-s)X_2$. Hence $(S,X)\in {\mathcal C}(t)$.
\end{proof}




\begin{claim}[Initial curvature lies in the pinching region]
\label{claim:pinching-initial-condition}
\uses{claim:pinching-region-convex}
${\mathcal T}(x,0)\in Z(x,0)$ for all $x\in M$.
\end{claim}

\begin{proof}
 Note that by the
hypothesis of the theorem we have
$$\nu(x,0)+\mu(x,0)+\lambda(x,0)\geq -3$$ so $(S(x,0),X(x,0))\in
{\mathcal C}(0)$ for all $x\in M$. On the other hand, if $0<
X(x,0)$, then since $X(x,0)\le 1$ we have $S(x,0)\geq-3X(x,0)\geq
X(\mathrm{log}X-3)$. This completes the proof that ${\mathcal
T}(x,0)\in {\mathcal C}(0)$ for all $x\in M$.
\end{proof}


\begin{claim}[The vector field preserves the pinching family of convex sets]
\label{claim:pinching-vector-field-preserves-family}
\uses{claim:pinching-region-convex}
The vector field $\psi({\mathcal T})={\mathcal T}^2+{\mathcal
T}^\sharp$ preserves the family ${\mathcal Z}(t)$ of convex sets.
\end{claim}

\begin{proof}
\uses{claim:pinching-auxiliary-quantity-nonnegative}
Fix $x\in M$ and suppose that we have an integral curve $\gamma(t),\
t_0\le t\le T$, for $\psi$ with $\gamma(t_0)\in Z(x,t_0)$. We wish
to show that $\gamma(t)\in Z(x,t)$ for all $t\in [t_0,T]$. The
function $S(t)=S(\gamma(t))$ satisfies
$$\frac{dS}{dt}=\lambda^2+\mu^2+\nu^2+\lambda\mu+\lambda\nu+\mu\nu=\frac{1}{2}\left((\lambda+\mu)^2+
(\lambda+\nu)^2+(\mu+\nu)^2\right).$$ By Cauchy-Schwarz we have
$$(\lambda+\mu)^2+
(\lambda+\nu)^2+(\mu+\nu)^2)\ge \frac{4S^2}{3}\ge \frac{2S^2}{3}.$$ Since
$\gamma(t_0)\in Z(x,t_0)$ we have $S(t_0)\ge -3/(1+t_0)$. It then follows that
\begin{equation}\label{Sineq}
S(t)\ge -3/(1+t)\ \ \ \text{for all}\ \ \  t\ge t_0.
\end{equation}


Now let us consider the evolution of $X(t)=X(\gamma(t))$. Assume
that we are at a point $t$ for which $X(t)>0$. For this computation
we set $Y=-\mu$.
\begin{align*}
\frac{dX}{dt}&=-\frac{d\nu}{dt}=-\nu^2-\mu\lambda =-X^2+Y\lambda,\\
\frac{dS}{dt}&=\frac{d(\nu+\mu+\lambda)}{dt}
=\nu^2+\mu^2+\lambda^2+\mu\lambda+\nu\lambda+\nu\mu\\
&=X^2+Y^2+\lambda^2+XY-\lambda(X+Y).
\end{align*}
Putting this together yields
\begin{equation}\label{dX}
X\frac{dS}{dt}-(S+X)\frac{dX}{dt} =X^3+I, \end{equation}
 where
$I=XY^2+\lambda Y(Y-X)+\lambda^2(X-Y)$.
\begin{claim}[Non-negativity of the auxiliary quantity I]
\label{claim:pinching-auxiliary-quantity-nonnegative}  $I\geq 0$.
\end{claim}

\begin{proof}

 First we consider the case when
 $Y\leq 0$. This means that  $\mu\geq 0$ and
 hence that $\lambda\geq0$. Since by definition $X\ge 0$, we have $X\ge  Y$.
This immediately gives  $I\geq 0$. Now let us consider the case when
$Y>0$ which means that $\nu\le \mu<0$. In this case, we have
\[ I=Y^3+(X-Y)(\lambda^2-\lambda Y+Y^2)>0 \]
since $X\geq Y$ and $\lambda^2-\lambda
Y+Y^2=(\lambda-\frac{Y}{2})^2+\frac{3Y^2}{4}>0$.
\end{proof}

The above claim and Equation~(\ref{dX}) immediately imply that
\begin{equation}\label{dX1}
 X\frac{dS}{dt}-(S+X)\frac{dX}{dt}\geq
X^3. \end{equation}
 Set $W=\frac{S}{X}-\mathrm{log}\,X$, then rewriting
 Equation~(\ref{dX1})
 in terms of $W$ gives
\begin{equation}\label{dW}
\frac{dW}{dt}\geq X.
\end{equation}

Now suppose that $\gamma(t)\not\in Z(x,t)$ for some $t\in [t_0,T]$.
Let $t_1<T$ be maximal subject to the condition that $\gamma(t)\in
Z(x,t)$ for all $t_0\le t\le t_1$. Of course, $\gamma(t_1)\in
\partial Z(x,t_1)$ which implies that $(S(t_1),X(t_1))\in \partial
{\mathcal C}(t_1)$. There are two possibilities: either
$S(t_1)=-3/(1+t_1)$ and $X(t_1)<\xi(t_1)$ or $X(t_1)\ge
\xi(t_1)>1/(1+t_1)$  and $S(t_1)=f_{t_1}(X(t_1))$. But
Equation~(\ref{Sineq}) implies that $S(t)\ge -3/(1+t)$ for all $t$.
Hence, if the first case holds then $\gamma(t)\in Z(x,t)$ for $t$ in
some interval $[t_0,t_1']$ with $t_1'>t_1$. This contradicts the
maximality of $t_1$. Thus, it must be the case that $X(t_1)\ge
\xi(t_1)$. But then $X(t)> \frac{1}{1+t}$ for all $t$ sufficiently
close to $t_1$. Hence, by Equation~(\ref{dW}) we have
\[ \frac{dW}{dt}(t)\geq X(t)>\frac{1}{1+t}, \]
for all $t$ sufficiently close to $t_1$. Also, since
$S(t_1)=f_{t_1}(X(t_1))$, we have $W(t_1)=(\mathrm{log}(1+t_1)-3)$. It
follows immediately that $W(t)\ge (\mathrm{log}(1+t)-3)$ for all
$t>t_1$ sufficiently close to $t_1$. This proves that $S(t)\ge
f_t(X(t))$ for all $t\ge t_1$ sufficiently close to $t_1$, again
contradicting the maximality of $t_1$.

This contradiction proves that $\psi$ preserves the family
${\mathcal Z}(t)$.
\end{proof}

By Theorem~\ref{varyingZ}, the previous three claims imply that ${\mathcal
T}(x,t)\in {\mathcal Z}(t)$ for all $x\in M$ and all $t\in [0,T)$. That is to
say, $S(x,t)\ge -3/(1+t)$ and $S(x,t)\ge f_t(X(x,t))$ whenever $X(x,t)\ge
\xi(t)$. For $X\in [1/(1+t),\xi(t)]$ we have $f_t(X)\le -3/(1+t)$, and thus in
fact $S(x,t)\ge f_t(X(x,t))$ as long as $X(x,t)\ge 1/(1+t)$. On the other hand,
if $0< X(x,t)\le 1/(1+t)$ then $f_t(X(x,t))< -3X(x,t)\le S(x,t)$. On the other
hand, since $X(x,t)$ is the negative of the smallest eigenvalue of ${\mathcal
T}(x,t)$ and $S(x,t)$ is the trace of this matrix, we have $S(x,t)\ge
-3X(x,t)$. Thus, $S(x,t)\ge f_t(X(x,t))$ in this case as well. This completes
the proof of Theorem~\ref{pinch}.
\end{proof}

Actually, the proof establishes a stronger result which we shall
need.

\begin{theorem}[Pinching toward positive curvature persists from time a]
\label{thm:pinching-persists-from-time-a}
\uses{thm:pinching-toward-positive-curvature}\label{pincha}
Fix $a\ge 0$. Let $(M,g(t)),\ a\le t<T$, be a Ricci flow with $M$ a
compact $3$-manifold. Suppose the eigenvalues of $\Rm(x,t)$ are
$\lambda(x,t)\ge \mu(x,t)\ge \nu(x,t)$ and set $X(x,t)=\mathrm{max}(-\nu(x,t),0)$. Assume that for every $x\in M$ we have
$$R(x,a) \ge  \frac{-6}{4a+1}$$
and $$ R(x,a) \geq 2X(x,a)\left(\mathrm{log}X(x,a)+\mathrm{log}(1+a)-3\right),$$ where the second inequality holds whenever
$X(x,a)>0$. Then for all $a\le t<T$ we have:
\begin{eqnarray}\label{Rlower}
R(x,t) & \ge & \frac{-6}{4t+1} \\\label{RXineq} R(x,t) & \geq &
2X(x,t)\left(\mathrm{log}X(x,t)+\mathrm{log}(1+t)-3\right),
\end{eqnarray}
whenever  $X(x,t)> 0$.
\end{theorem}

Once again it is natural to interpret the right-hand side of the
inequalities relating $R$ and $X$  to be zero when $X(x,t)=0$. With
this convention the result remains true even when $X(x,t)=0$.


\begin{corollary}[Curvature bound from a scalar-curvature bound under pinching]
\label{cor:curvature-bound-from-scalar-pinching}
\uses{thm:pinching-persists-from-time-a}
Fix $a\ge 0$. Suppose that $(M,g(t)),\ a\le t<T$, is a Ricci flow with $M$ a
compact $3$-manifold, and suppose that the two hypotheses of the previous
theorem hold. Then there is a continuous function $\phi$ such that for all
$R_0<\infty$, if $R(x,t)\le R_0$ then $|\Rm(x,t)|\le \phi(R_0)$.
\end{corollary}


\begin{proof}
\uses{thm:pinching-persists-from-time-a}
Fix $R_0\ge e^4$ sufficiently large, and suppose that $R(x,t)\le
R_0$. If $X(x,t)=0$, then $|\Rm(x,t)|\le R(x,t)/2$. If
$X(x,t)>0$, then by Theorem~\ref{pincha} it is bounded by $R_0$.
Thus, $\lambda(x,t)\le 3R_0$. Thus, we have an upper bound on
$\lambda(x,t)$ and a lower bound on $\nu(x,t)$ in terms of $R_0$.
\end{proof}

This theorem leads to a definition.

\begin{definition}[Curvature pinched toward positive for a generalized Ricci flow]
\label{def:curvature-pinched-toward-positive}
\uses{def:generalized-ricci-flow,thm:pinching-toward-positive-curvature}\label{pinchd}
Let $({\mathcal M},G)$ be a generalized Ricci flow whose domain of definition
is contained in $[0,\infty)$. Then we say that $({\mathcal M},G)$ has {\em
curvature pinched toward positive} if
for every $x\in {\mathcal M}$ the following two conditions hold:
\begin{enumerate}
\item $$R(x)\ge \frac{-6}{4{\bf t}(x)+1}$$
\item $$R(x)\geq 2X(x)\left(\mathrm{log}X(x)+\mathrm{log}(1+{\bf t}(x))-3\right),
$$ whenever
$0<X(x)$
\end{enumerate}
where, as in the statement of Theorem~\ref{pinch}, $X(x)$ is the
maximum of zero and the negative of the smallest eigenvalue of $\Rm(x)$.
\end{definition}

The content of Theorem~\ref{pincha} is that if $(M,g(t)),\ 0\le a\le
t<T$, is a Ricci flow with $M$ a compact $3$-manifold and if the
curvature of $(M,g(a))$ is pinched toward positive, then the same is
true for the entire flow.


\subsection{Application of the pinching result}


As an application of this pinching toward positive curvature result
we establish a strengthening of Theorem~\ref{sphsf}.


\begin{theorem}[Compact shrinking solitons are round]
\label{thm:shrinking-soliton-round}
Let $(M,g)$ be a compact $3$-dimensional shrinking soliton, i.e.,
there is a Ricci flow $(M,g(t)),\ 0\le t<T$, so that for each $t\in
[0,T)$ there is a constant $c(t)$ with  $\mathrm{lim}_{t\rightarrow
T}c(t)=0$ and with the property that there is an isometry from
$(M,g(t))$ to $(M,c(t)g)$. Then $(M,g)$ is round.
\end{theorem}

\begin{proof}
\uses{thm:nonnegative-ricci-flow-becomes-round,thm:pinching-toward-positive-curvature,cor:nonnegative-curvature-operator-preserved}
By rescaling we can assume that for all $x\in M$ all the eigenvalues of $\Rm(x,0)$ have absolute value $\le 1$.  This implies that $(M,g(0))$ satisfies
the hypothesis of Theorem~\ref{pinch}. Our first goal is to show that $\Rm(x,0)\ge 0$ for all $x\in M$. Suppose that this is not true; then there is a
point $x$ with $X(x,0)>0$. Consider $A=X(x,0)/R(x,0)$. For each $t<T$ let
$x_t\in M$ be the image of $x$ under the isometry from $(M,g(0))$ to
$(M,c(t)g(t))$.
 Then $X(x_t,t)=c^{-1}(t)X(x,0)$ and  $X(x_t,t)/R(x_t,t)=A$. Since
$c(t)$ tends to $0$ as $t$ approaches $T$, this contradicts
Theorem~\ref{pinch}. Now, according to Theorem~\ref{rm>0} either all
$(M,g(t))$ are flat or $\Rm(x,t)>0$ for all $(x,t)\in M\times
(0,T)$. But if the $(M,g(t))$ are all flat, then the flow is trivial
and hence the diameters of the $(M,g(t))$ do not go to zero as $t$
approaches $T$, contradicting the hypothesis. Hence, $\Rm(x,t)>0$ for all $(x,t)\in M\times (0,T)$. According to
Theorem~\ref{flowtoround} this means that as the singularity
develops the metrics are converging to round. By the shrinking
soliton hypothesis, this implies that all the metrics $(M,g(t)),\
0<t<T$, are in fact round. Of course, it then follows that $(M,g)$
is round.
\end{proof}


The following more general result was first given by T. Ivey
\cite{Ivey}.

\begin{theorem}[Compact three-dimensional Ricci solitons are Einstein]
\label{thm:ricci-soliton-einstein-dimension-three}
\uses{def:einstein-manifold}\label{ivey}
Any 3-dimensional compact Ricci soliton
$g_0$ is Einstein.
\end{theorem}

Since we do not need this result, we do not include a proof.



\subsection{The Harnack inequality}

The last consequence of the maximum principle that we need is
Hamilton's version of the  Harnack inequality for Ricci
flows, see Theorem 1.1
and Corollary 1.2 of \cite{Hamiltonharnack}.

\begin{theorem}[Hamilton's Harnack inequality for Ricci flow]
\label{thm:hamilton-harnack-inequality}
\uses{def:ricci-curvature,def:curvature-operator}\label{Harnack}
Suppose that $(M,g(t))$ is a Ricci flow defined for $(T_0,T_1)$ with
$(M,g(t))$ a complete manifold of non-negative curvature operator
with bounded curvature  for each $t\in (T_0,T_1)$. Then for any
time-dependent vector field $\chi(x,t)$ on $M$ we have:
$$\frac{\partial R(x,t)}{\partial t}+\frac{R(x,t)}{t-T_0}+2\langle
\chi(x,t),\nabla R(x,t)\rangle +2\Ric(x,t)(\chi(x,t),\chi(x,t))\ge 0.$$ In particular, we have
$$\frac{\partial R(x,t)}{\partial t}+\frac{R(x,t)}{t-T_0}\ge 0.$$
\end{theorem}



\begin{remark}[Harnack inequality special case with zero vector field]
\label{rem:harnack-inequality-zero-field-case}
\uses{thm:hamilton-harnack-inequality}
Notice that the second result follows from the first by taking
$\chi=0$.
\end{remark}

\begin{corollary}[Scalar curvature is non-decreasing for ancient flows]
\label{cor:ancient-flow-scalar-curvature-nondecreasing}
\uses{thm:hamilton-harnack-inequality}\label{posderiv}
If $(M,g(t))$ is a Ricci flow defined for $-\infty<t\le 0$ with
$(M,g(t))$ a complete manifold of bounded, non-negative curvature
operator for each $t$, then
$$\frac{\partial R(x,t)}{\partial t}\ge 0.$$
\end{corollary}

\begin{proof}
\uses{thm:hamilton-harnack-inequality}
Apply the above theorem with $\chi(x,t)=0$ for all $(x,t)$ and for a
sequence of $T_0\rightarrow -\infty$.
\end{proof}


The above is the differential form of Hamilton's Harnack inequality.
There is also the integrated version, also due to Hamilton; see
Corollary 1.3 of \cite{Hamiltonharnack}.

\begin{theorem}[Integrated Harnack inequality]
\label{thm:harnack-inequality-integrated}
\uses{def:geodesic}\label{Harnack2}
Suppose that $(M,g(t))$ is a Ricci flow defined for $t_1\le t\le
t_2$ with $(M,g(t))$ a complete manifold of non-negative, bounded
curvature operator for all $t\in [t_1,t_2]$. Let $x_1$ and $x_2$ be
two points of $M$.  Then $$\mathrm{log}\left(\frac{R(x_2,t_2)}{R(x_1,t_1)}\right)\ge
-\frac{1}{2}\frac{d^2_{t_1}(x_2,x_1)}{(t_2-t_1)}.$$
\end{theorem}


\begin{proof}
\uses{thm:hamilton-harnack-inequality}
Apply the differential form of the Harnack inequality to
$\chi=-\nabla(\mathrm{log} R)/2=-\nabla R/2R$, and divide by $R$. The
result is
$$R^{-1}(\partial R/\partial t)-|\nabla (\mathrm{log} R)|^2+\frac{\Ric(\nabla(\mathrm{log} R),\nabla(\mathrm{log} R))}{2R}\ge 0.$$ Since $\Ric(A,A)/R\le |A|^2$, it follows
that
$$\frac{\partial}{\partial t}(\mathrm{log} R)-\frac{|\nabla(\mathrm{log} R)|^2}{2}\ge 0.$$

Let $d$ be the $g(t_1)$-distance from $x_1$ to $x_2$ and let
 $\gamma\colon [t_1,t_2]\to M$ be a $g(t_1)$-geodesic from $x_1$
 to $x_2$, parameterized at speed $d/(t_2-t_1)$. Then
 let $\mu(t)=(\gamma(t),t)$ be a path in space-time.
 We compute
 \begin{eqnarray*}
 \mathrm{log}\left(\frac{R(x_2,t_2)}{R(x_1,t_1)}\right) & = & \int_{t_1}^{t_2}
 \frac{d}{dt}\mathrm{log}(R(\mu(t))dt \\
 & = & \int_{t_1}^{t_2}\frac{\frac{\partial R}{\partial t}(\mu(t))}{R(\mu(t))}+\langle\nabla(\mathrm{log} R)(\mu(t)),\frac{d\mu}{dt}(\mu(t))dt \\
 & \ge & \int_{t_1}^{t_2}\frac{1}{2}|\nabla (\mathrm{log} R)(\mu(t))|^2-|\nabla (\mathrm{log} R)(\mu(t))|\cdot \left|\frac{d\gamma}{dt}\right|dt \\
& \ge  & -\frac{1}{2}\int_{t_1}^{t_2}\left|\frac
{d\gamma}{dt}\right|^2dt,\end{eqnarray*} where the last inequality
comes form completing the square.
 Since $\Ric(x,t)\ge 0$,
$|d\gamma/dt|_{g(t)}\le |d\gamma/dt|_{g(t_1)}$, thus
$$\mathrm{log}\left(\frac{R(x_2,t_2)}{R(x_1,t_1)}\right)\ge
-\frac{1}{2}\int_{t_1}^{t_2}\left|\frac{d\gamma}{dt}\right|^2_{g(t_1)}dt.$$
Since $\gamma$ is a $g(t_1)$-geodesic, this latter integral is
$$-\frac {1}{2}\frac{d_{g(t_1)}^2(x_1,x_2)}{t_2-t_1}.$$
\end{proof}
